\section{Holomorphic transfer pencils and symmetry}
\label{sec:related}

Symbolic zeta determinants and holomorphic transfer families are classical
\citep{bowenlanford1970,ruelle1976}.  Group twists organize periodic words by
their cocycle monodromy \citep{parrypollicott1990}, while identity-word zeta
coefficients for matrices over free-group algebras provide the algebraic
precedent for the positive-cone trace \citep{kasselreutenauer2014}.

Flip systems and reversal signatures already provide nontrivial symbolic
matrix formalisms for time reversal \citep{kimleepark2003,ryu2023}.  Their
reversal-fixed counts are different from the tautological channel exchange
used here.  Classical regularized determinants supply the $\det_q$ expansion
\citep{simon1977}; they do not turn block exchange into arithmetic content.

The present question is narrower than a general transfer-operator functional
equation.  We impose $s\leftrightarrow1-s$ directly through a two-layer
symbolic grammar and ask whether one holomorphic determinant can both retain
the tensor-prime ledger and move vertically.  This differs from
self-adjointization: no conjugation or adjoint appears in the definition.

Formal-language zeta theory shows that cyclic languages can admit
finite-dimensional virtual-character realizations
\citep{berstelreutenauer1990}; exterior automata give alternating
determinants for sofic systems \citep{beal1995}.  These results warn that
finite channels and signs can change a trace presentation without changing
the language's exponent ledger.  Our rigidity theorem makes that warning
explicit for reflected tensor-prime monomials.

We found no primary source proving the precise balance dichotomy used here:
a finite-channel reflected positive-cone double over entropy-derived
tensor-prime atoms is either vertically sterile on pure identity-visible
words or admits mixed arithmetic masses.  We claim this scoped theorem and
its source-locked synthesis, not novelty for holomorphic transfer operators,
twisted zeta functions, or formal-language determinants.
